\section{Introduction}
\label{sec:intro}

Many datasets are made of relations among more than two things. A medical claim involves a patient, a provider, a payer and a list of procedures. A paper has authors in a fixed order. A reaction has reactants, products and a catalyst, each in its own role. A hypergraph models each such relation as one \emph{hyperedge} that contains any number of members, where a member can play a role, carry a weight or hold only for a period of time. If hyperedges are objects in their own right, a hyperedge can also be a member of another one. Statements about statements, provenance and evidence are expressed this way~\cite{wikidata,battiston}.

Graph databases store binary edges. The usual way to store an $n$-ary relation in one is reification: a node stands for the relation and $n$ binary edges connect it to its members~\cite{angles}. Queries still work, but the storage engine now sees $n+1$ unrelated objects. It cannot say how many members a relation has without visiting all $n$ edges or compare two relations without reading both, and a change to a membership becomes $n$ separate edge updates that a concurrent writer cannot merge as a unit.

Several systems model hyperedges directly. HyperGraphDB~\cite{hgdb} stores each link as an array of target handles in Berkeley DB Java Edition~\cite{bdb}, with a separate index from every atom to the links that point at it. AtomSpace~\cite{opencog} and TypeDB~\cite{typeql} make $n$-ary relations first class in their data models, and parallel frameworks~\cite{hygra,mesh,hyperx} analyse hypergraphs held in memory. What is still missing is a storage engine whose physical structures are designed around hyperedges. Such an engine should make small and very large memberships equally cheap to read and change, answer membership, cardinality and overlap questions without scanning, let writers commit concurrently while readers see stable snapshots and past versions, and recover from a crash to a state that some prefix of the commits produced.

This paper describes HStore, an embedded storage engine built to those requirements, and compares it with HyperGraphDB. Four ideas shape it.

\paragraph{One tree for every structure.} Every structure HStore keeps is a copy-on-write B+tree over 64-bit keys, including each hyperedge's membership and the reverse index from an atom to the hyperedges that contain it. Each child reference in the tree carries the entry count and a summary of its subtree: the key range, the weight and time ranges, a role mask and a fingerprint that depends only on the subtree's contents. Counts give rank and select in logarithmic time. Summaries let a query skip subtrees it cannot match, and let two trees be compared, intersected or diffed in time that grows with how much they differ rather than with their size. The catalog is a tree of trees: a hyperedge's record holds the root of its membership tree, and holds the whole tree when it is small, as almost all are.

\paragraph{Commits log keys, not pages.} A copy-on-write commit normally writes every node from each changed leaf up to the root. An HStore commit logs the keys it changed and publishes a generation whose new nodes are still in memory. A background checkpoint writes each of those nodes once, however many commits changed it. Recovery replays the logged keys from the last checkpoint in commit order and stops at the first commit whose pages did not reach the disk, so the recovered state is always a prefix of the commit order.

\paragraph{Pages move without rewriting trees.} Trees refer to pages by permanent page numbers, and a memory-mapped page directory maps each number to the page's current place in an append-only segment file. Compaction copies a live page byte for byte and changes one directory entry. No tree, generation or branch is rewritten, so compaction does not conflict with commits and keeps old generations readable. Reference counts, updated only between consecutive checkpointed generations, tell compaction which pages are live.

\paragraph{Leaves decode one record at a time.} Records in catalog and index leaves are prefixed by their length. A leaf read from disk keeps its bytes and decodes a record only when the record is first read, charging the decoded size to the node cache. When the data is larger than the cache, a lookup then decodes one record instead of a whole page.

HStore is written in Java~25 and consists of the storage engine, a database layer with schemas, properties and a query language, and a server. We compare it with HyperGraphDB on a synthetic hypergraph with skewed degrees, using \M{workloadCount} workloads that cover ingest, reads, per-operation latency, reads during writes, deletes, space reclamation and crash recovery. We run them with asynchronous and synchronous commits, and with a dataset larger than both systems' caches. HStore is ahead on every workload whose rate the harness does not fix, with one exception: the median latency of a single commit under asynchronous durability, where HyperGraphDB is faster. HStore writes \M{writtens1asyncwriteupdatehs} per update where HyperGraphDB writes \M{writtens1asyncwriteupdatehg}, and its files after the load are \M{ratios1asyncdisk} to \M{ratios4asyncdisk} times smaller. It keeps more of the Java heap than HyperGraphDB does.

\Cref{sec:model} defines the data model. \Cref{sec:overview} gives an overview of the engine, and \cref{sec:layout,sec:transactions,sec:space} describe the storage layout, transactions and recovery, and space reclamation. \Cref{sec:implementation} covers the implementation and its testing, and \cref{sec:evaluation} the evaluation. \Cref{sec:related} discusses related work, and \cref{sec:conclusion} concludes with the limitations we know of.
